\BeleskeID{05-s023}
% element: 05-s023:n01
% element: 05-s023:d01
Relacije sa oznakom $\equiv$ predstavljaju jednakost sačuvanih $n$ bitova, a ne nužno običnu jednakost svih celobrojnih zbirova. Za $a<b$ zbir sa komplementom drugog operanda daje zapis negativnog rezultata bez izlaznog prenosa. Za $a>b$ pojavljuje se dodatni $2^n$, odnosno izlazni prenos; odbacivanjem tog prenosa ostaje pozitivan rezultat.

Brojevi 5 i 12 u ovom primeru imaju pet bitova kako bi oba bila pozitivna u drugom komplementu. Zapis negativnog broja $-12$ je $32-12=20=10100_2$. Zbir sa 5 je $25=11001_2$. Za čitanje apsolutne vrednosti rezultata invertiramo bite: $00110$, pa dodamo jedan: $00111=7$. Znak najvišeg bita zato znači da je rezultat $-7$.
